\chapter{补充复查表格}
\label{cha:appendix_review_tables}

本附录给出实验复查时可使用的记录字段。附录表格不替代正文实验结果，而是用于说明困难样本和部署一致性复查应保留哪些信息。

\section{困难样本复查字段}

表 \ref{tab:appendix_hardcase_template} 用于记录困难样本的观察结果。该表仅保留论文复查所需字段，不展开本地文件组织细节，也不把样本整理为新的人工分组实验。

\begin{table}[htbp]
    \centering
    \caption{困难样本复查字段}
    \label{tab:appendix_hardcase_template}
    \begin{tabular}{p{2.5cm} p{3.2cm} p{3.2cm} p{3.0cm}}
        \toprule
        字段 & 记录内容 & 检测表现 & 备注 \\
        \midrule
        样本编号 & 图像或视频帧标识 & 是否有漏检、误检或重复框 & 对应日志或截图路径 \\
        场景描述 & 文字描述可见干扰 & 类别分数和边框变化 & 不作为统计分组 \\
        复查结论 & 关联阈值、NMS 或后处理设置 & 是否需要进入正文讨论 & 对应表格或图编号 \\
        \bottomrule
    \end{tabular}
\end{table}

\section{部署一致性复查字段}

表 \ref{tab:appendix_deploy_check} 给出了部署一致性复查时应记录的字段。训练阶段、ONNX 导出阶段、OM 编译阶段和板端部署阶段的逐样本对齐实验可沿用该表记录类别、边框、阈值和运行状态。

\begin{table}[htbp]
    \centering
    \caption{部署一致性复查字段}
    \label{tab:appendix_deploy_check}
    \begin{tabular}{p{2.8cm} p{4.0cm} p{4.2cm}}
        \toprule
        复查项 & 记录内容 & 判断依据 \\
        \midrule
        类别顺序 & person / ebike 是否一致 & 类别配置与输出解释一致 \\
        输入预处理 & 尺寸、缩放、颜色通道 & 输入契约一致 \\
        输出张量 & 坐标格式、分数含义 & 输出契约一致 \\
        后处理参数 & conf、NMS 阈值 & 阈值行为可比较 \\
        运行状态 & 成功样本数、失败样本数、平均时延 & 板端结果可复查 \\
        \bottomrule
    \end{tabular}
\end{table}
